% Seção: Plano de desenvolvimento
% Dono: Henrique · Orçamento: 0,5 página(s) · docs/fontes/plano-de-projeto.md §14.1
% Conteúdo: Método, divisão e calendário (Etapa 1)
% Fontes: docs/fontes/plano-de-projeto.md §1, §6 e §7; plano.md (finalização de 20/11).
\section{Plano de desenvolvimento}\label{sec:plano}

O trabalho é dividido por \emph{feature}: cada uma responde a uma das seis perguntas e tem um
dono de ponta a ponta (algoritmo, interface, testes, medição, guia, seção do relatório e vídeo).
Cada integrante também cuida de uma frente da fundação: núcleo do grafo, dados, modelagem,
interface, repositório e avaliação. As interfaces entre as partes são publicadas como contratos
antes do código que depende delas. A \texttt{main} só recebe código por \emph{pull request}
aprovado por outro integrante, e um gancho de \emph{pre-commit} barra qualquer biblioteca de
grafo, e também o \texttt{heapq}, em \texttt{src/}; a NetworkX~\cite{hagberg2008} aparece só
nos testes e nos scripts de medição, para comparar resultados. O calendário teve data fixa e
escopo variável: fundação até 03/11, essenciais até 11/11, desejáveis até 15/11 e entrega até
22/11. Em 20/11 o grupo suspendeu o congelamento de código de 15/11 para fechar as extensões
desejáveis e opcionais (Seção~\ref{sec:extensoes}).
